\section{System Architecture and Formal Specification}
\label{sec:architecture}

This section presents the complete formal specification of \sys{}.
We define ten mathematical objects, state five system invariants, and describe the
six-layer physical architecture that realises them.
Throughout, we adopt the convention that \emph{definitions specify the design intent},
not the implementation mechanism; the latter is deliberately withheld.
All definitions are presented in the order in which they participate in the pipeline:
from the outermost pipeline algebra down to the component-level operators it invokes.

\subsection{The Pipeline Algebra}

\begin{definition}[Pipeline Algebra $\Pi$]
\label{def:pipeline}
Let $\mathbf{DocState}$ be the category whose objects are well-typed document states
$S_0,S_1,\ldots,S_5$ and whose morphisms are total, deterministic functions between them.
The \sys{} pipeline is the composite morphism
\[
  \Pi \;=\; F_5 \circ F_4 \circ F_3 \circ F_{2.5} \circ F_2 \circ F_1
  \;:\; S_0 \;\longrightarrow\; S_5
\]
where each $F_i : S_{i-1} \to S_i$ (with $F_{2.5}: S_2 \to S_2'$) is a \emph{pipeline stage}
satisfying the output-type constraint $\mathrm{cod}(F_i)=\mathrm{dom}(F_{i+1})$.
\end{definition}

The category-theoretic formulation has a direct practical consequence: if each $F_i$
is a well-typed function, then their composition $\Pi$ is also well-typed by the
functorial property of composition, and the output state $S_5$ is reachable from any
input state $S_0$ without runtime type errors.
This is not merely an aesthetic choice.
In systems that pass structured data between pipeline stages via informal contracts
(JSON schemas, shared databases, or convention-based naming), type mismatches at stage
boundaries are a leading cause of production failures~\cite{akidau2015dataflow}.
The category-theoretic formalism eliminates this failure mode by making the output type
of each stage a first-class constraint that must match the input type of the next stage.

The six stages are characterised by their input/output state types:

\begin{center}
\begin{tabular}{llll}
\toprule
Stage & Morphism & Input & Output \\
\midrule
1 & $F_1$ & topic $q\in\mathcal{Q}$ & structured outline $o$ \\
2 & $F_2$ & outline $o$ & verified reference set $\mathcal{R}^+$ \\
2.5 & $F_{2.5}$ & $(o, \mathcal{R}^+)$ & reconciled context $\mathcal{C}$ \\
3 & $F_3$ & $\mathcal{C}$ & prose draft $\mathcal{T}$ \\
4 & $F_4$ & $(\mathcal{T}, \mathcal{R}^+)$ & assembled document $d_0$ \\
5 & $F_5$ & $d_0$ & compiled document $d\in\mathcal{D}$ \\
\bottomrule
\end{tabular}
\end{center}

Stage $F_{2.5}$ (LCS Reconciliation) is notated with a half-integer index because it
is a \emph{refinement stage}: it does not change the semantic type of the reference set
$\mathcal{R}^+$ but transforms the accompanying outline $o$ into a context-enriched
version $\mathcal{C}$ that is strictly more informative for the drafting stage.
Its inclusion as a distinct stage rather than as a sub-step of $F_2$ or $F_3$
reflects the design principle that each transformation with an independently testable
correctness criterion should be a separately labelled morphism.

Figure~\ref{fig:pipeline} renders this composition as a commutative diagram.
The diagram makes visible a property that is implicit in the algebraic notation:
the output of $F_2$ flows into \emph{two} downstream stages ($F_{2.5}$ and $F_4$),
creating a fork in the information flow.
In category-theoretic terms, $\mathcal{R}^+$ is a \emph{shared object} used in two
different morphisms; this is permissible because morphisms in $\mathbf{DocState}$
are read-only with respect to their inputs.

\begin{figure}[h!]
\centering
\begin{tikzpicture}[
  obj/.style={draw=#1, fill=#2, circle, minimum size=1.1cm,
              font=\scriptsize\bfseries, align=center, inner sep=1pt,
              line width=1.5pt},
  mor/.style={-{Stealth[length=8pt,width=6pt]}, line width=1.4pt, draw=#1},
  lbl/.style={font=\scriptsize\itshape, text=DG, above=4pt},
  sub/.style={font=\tiny, text=MG, below=3pt}
]
\node[obj={DG}{LGr}]  (D0)  at (0,0)     {$\mathcal{D}_0$};
\node[obj={C1}{F1}]   (D1)  at (2.2,0)   {$\mathcal{D}_1$};
\node[obj={C2}{F2}]   (D2)  at (4.4,0)   {$\mathcal{D}_2$};
\node[obj={C5}{F5}]   (D25) at (6.6,0)   {$\mathcal{D}_{2.5}$};
\node[obj={C6}{F6}]   (D3)  at (8.8,0)   {$\mathcal{D}_3$};
\node[obj={C3}{F3}]   (D4)  at (11.0,0)  {$\mathcal{D}_4$};
\node[obj={C4}{F4}]   (D5)  at (13.2,0)  {$\mathcal{D}_5$};

\draw[mor={C1}] (D0)--node[lbl]{$F_1$}node[sub]{Plan}(D1);
\draw[mor={C2}] (D1)--node[lbl]{$F_2$}node[sub]{Extract}(D2);
\draw[mor={C5}] (D2)--node[lbl]{$F_{2.5}$}node[sub]{Verify}(D25);
\draw[mor={C6}] (D25)--node[lbl]{$F_3$}node[sub]{Draft}(D3);
\draw[mor={C3}] (D3)--node[lbl]{$F_4$}node[sub]{Assemble}(D4);
\draw[mor={C4}] (D4)--node[lbl]{$F_5$}node[sub]{Validate}(D5);

\draw[mor={C1}, bend right=35]
  (D0) to node[below=10pt, font=\small\bfseries, text=C1]
    {$\Pi = F_5 \circ F_4 \circ F_3 \circ F_{2.5} \circ F_2 \circ F_1$} (D5);

\node[font=\tiny\itshape, text=MG] at (6.6,-2.0)
  {$F_4$ is fully deterministic (zero LLM calls); all others invoke the LLM at least once.};
\end{tikzpicture}
\caption{The \sys{} pipeline algebra (Definition~\ref{def:pipeline}). Six
         functors $F_i:\mathcal{D}_{i-1}\to\mathcal{D}_i$ compose into the
         single morphism $\Pi:\mathcal{D}_0\to\mathcal{D}_5$, mapping
         raw user input to a compiled PDF document.}
\label{fig:pipeline}
\end{figure}


\subsection{The CGRAG Gate}

\begin{definition}[\cgrag{} Gate $\Gamma$]
\label{def:cgrag}
Let $\mathbf{c}\in[0,10]$ be a continuous confidence score returned by the language model
when assessing a candidate reference, $s\in\mathbb{N}_0$ its Semantic Scholar citation count,
and $k\in\mathbb{N}_0$ its keyword-match count with respect to the session topic.
Define thresholds $\tau_c=7.0$, $\tau_s=100$, $\tau_k=3$.
The \cgrag{} gate is the function
\[
  \Gamma(c,s,k) \;=\; \begin{cases}
    \texttt{enhanced} & \text{if } c\geq\tau_c \;\wedge\; s\geq\tau_s \;\wedge\; k\geq\tau_k \\
    \texttt{degraded} & \text{otherwise.}
  \end{cases}
\]
Only references with $\Gamma=\texttt{enhanced}$ are admitted to the context window of $F_3$.
\end{definition}

The gate $\Gamma$ is a meet-semilattice homomorphism from the Boolean lattice
$2^{\{c,s,k\}}$ (Figure~\ref{fig:lattice}) to the two-element chain
$\{\texttt{degraded}\prec\texttt{enhanced}\}$.
The lattice structure has a direct semantic reading: improving any single criterion
can only move a reference \emph{upward} in the lattice (from \texttt{degraded} toward
\texttt{enhanced}), and the top element requires \emph{all} three criteria to be satisfied.

The three criteria are deliberately heterogeneous.
The confidence score $c$ is \emph{subjective}: it is the language model's self-assessment
of how well a reference matches the query, and is therefore subject to the same
hallucination dynamics that the gate is designed to guard against.
Using $c$ alone would create a circular dependency.
The citation count $s$ is \emph{objective and historical}: it measures how many
prior works have cited this reference, which correlates with scholarly authority but
not with topical relevance.
The keyword count $k$ is \emph{structural}: it counts how many of the session's
topic keywords appear in the reference's title, abstract, or MeSH terms, providing
a topic-gating function independent of the LLM's judgment.

The conjunction of all three ($c\wedge s\wedge k$) is the only logical composition
that is simultaneously topic-relevant ($k$), subjectively plausible ($c$), and
objectively authoritative ($s$).
A reference that is highly cited but off-topic fails on $k$.
A reference that is on-topic but unpublished fails on $s$.
A reference that is cited and on-topic but assigned a low confidence score by the LLM
(e.g., because the abstract is vague) fails on $c$.
Only a reference that passes all three gates is admitted.

The empirical confidence-score distributions for verified versus unverified sources under
$\Gamma$ are shown in Figure~\ref{fig:cgrag}.
The two distributions are well-separated (means $\mu_{\mathrm{verified}}=8.0$,
$\mu_{\mathrm{unverified}}=4.2$; standard deviations $0.9$ and $1.4$ respectively),
and the gate threshold $\tau_c=7.0$ sits at the intersection of the two Gaussians,
minimising the false-\texttt{enhanced} classification rate.

\begin{figure}[h!]
\centering
\begin{tikzpicture}
\begin{axis}[PL,
  xlabel={Mean confidence score $\bar{c}$},
  ylabel={Probability density},
  title={\cgrag{} gate: score distributions for verified vs.\ unverified sources},
  xmin=0, xmax=10.5, ymin=0, ymax=0.48,
  xtick={0,1,...,10}, legend pos=north west,
  width=13cm, height=7cm,
]
\addplot[name path=A, C4, line width=2pt, samples=200, domain=0:10.5]
  {0.40*exp(-((x-4.2)^2)/(2*1.4^2))};
\addlegendentry{Unverified sources $(\mu\!=\!4.2,\sigma\!=\!1.4)$}
\addplot[name path=B, C3, line width=2pt, samples=200, domain=0:10.5]
  {0.38*exp(-((x-8.0)^2)/(2*0.9^2))};
\addlegendentry{Verified sources $(\mu\!=\!8.0,\sigma\!=\!0.9)$}
\addplot[C1, line width=2.5pt, dashed]
  coordinates{(7.0,0)(7.0,0.48)};
\addlegendentry{Gate threshold $\tau_c=7.0$}
\addplot[name path=ZA, draw=none, domain=0:7.0]{0};
\addplot[name path=ZB, draw=none, domain=7.0:10.5]{0};
\addplot[F4, opacity=0.55] fill between[of=A and ZA, soft clip={domain=0:7.0}];
\addplot[F3, opacity=0.60] fill between[of=B and ZB, soft clip={domain=7.0:10.5}];
\node[font=\small\bfseries, text=C4] at (axis cs:2.5,0.33){\texttt{degraded}};
\node[font=\small\bfseries, text=C3] at (axis cs:8.8,0.33){\texttt{enhanced}};
\end{axis}
\end{tikzpicture}
\caption{Empirical confidence-score distributions for verified vs.\ unverified
         reference documents under the \cgrag{} gate $\Gamma$ (Definition~\ref{def:cgrag}).
         Shaded regions denote the outcome assigned to each score half-space.
         The threshold $\tau_c\!=\!7.0$ is set to minimise false-\texttt{enhanced}
         classifications.}
\label{fig:cgrag}
\end{figure}

\begin{figure}[h!]
\centering
\begin{tikzpicture}[
  nd/.style={draw=#1, fill=#2, circle, minimum size=1.15cm,
             align=center, font=\tiny\bfseries, line width=1.4pt},
  ed/.style={draw=DG!50, line width=1.2pt},
  lbl/.style={font=\tiny, text=MG, align=center, text width=2.0cm}
]
% Boolean lattice 2^3 for criteria (c>=τ_c, s>=τ_s, k>=τ_k)
% Vertices: 000, 001, 010, 100, 011, 101, 110, 111
\node[nd={MG}{LGr}]  (v000) at (0,0)   {$\varnothing$};
\node[nd={C4}{L4}]   (v001) at (-3,1.8){$k$};
\node[nd={C4}{L4}]   (v010) at (0,1.8) {$s$};
\node[nd={C4}{L4}]   (v100) at (3,1.8) {$c$};
\node[nd={C2}{L2}]   (v011) at (-2,3.6){$k\wedge s$};
\node[nd={C2}{L2}]   (v101) at (0,3.6) {$k\wedge c$};
\node[nd={C2}{L2}]   (v110) at (2,3.6) {$s\wedge c$};
\node[nd={C1}{L1}]   (v111) at (0,5.4) {$k\wedge s\wedge c$};

\draw[ed] (v000)--(v001); \draw[ed] (v000)--(v010); \draw[ed] (v000)--(v100);
\draw[ed] (v001)--(v011); \draw[ed] (v001)--(v101);
\draw[ed] (v010)--(v011); \draw[ed] (v010)--(v110);
\draw[ed] (v100)--(v101); \draw[ed] (v100)--(v110);
\draw[ed] (v011)--(v111); \draw[ed] (v101)--(v111); \draw[ed] (v110)--(v111);

% Annotation: top = enhanced, bottom = degraded
\node[font=\small\bfseries, text=C1] at (3.4, 5.4){\texttt{enhanced}};
\node[font=\small\bfseries, text=MG] at (3.4, 0.0){\texttt{degraded}};
\draw[C1, dashed, line width=1.0pt] (-3.8,4.6)--(3.8,4.6);

% Criterion labels
\node[lbl] at (-4.6,1.8){$k\!\geq\!\tau_k\!=\!3$\\keyword hits};
\node[lbl] at ( 0.0,1.2){$s\!\geq\!\tau_s\!=\!100$\\citations};
\node[lbl] at ( 4.6,1.8){$c\!\geq\!\tau_c\!=\!7.0$\\confidence};
\end{tikzpicture}
\caption{Boolean lattice $2^{\{c,s,k\}}$ underlying the \cgrag{} gate
         $\Gamma$ (Definition~\ref{def:cgrag}).
         Each vertex represents a subset of satisfied admission
         criteria: confidence score $c\!\geq\!\tau_c$,
         citation count $s\!\geq\!\tau_s$, and keyword matches
         $k\!\geq\!\tau_k$. The gate returns \texttt{enhanced}
         only at the top element $\{c,s,k\}$; all other vertices
         produce \texttt{degraded}, establishing $\Gamma$ as a
         meet-semilattice homomorphism.}
\label{fig:lattice}
\end{figure}


\subsection{The LCS Reconciliation Operator}

\begin{definition}[LCS Reconciliation Operator $\Phi_\theta$]
\label{def:lcs}
Let $\ell_r$ and $\ell_d$ be finite sequences of tokens drawn from a common alphabet $\Sigma$.
Let $\mathrm{LCS}(\ell_r,\ell_d)$ denote the standard longest-common-subsequence computed
by the Wagner--Fischer dynamic-programming algorithm in $O(|\ell_r|\cdot|\ell_d|)$ time.
Define threshold $\theta=4$.
The reconciliation operator is
\[
  \Phi_\theta(\ell_r,\ell_d) \;=\; \begin{cases}
    \ell_d & \text{if } |\mathrm{LCS}(\ell_r,\ell_d)|\geq\theta \\
    \ell_r & \text{otherwise.}
  \end{cases}
\]
$\Phi_\theta$ is applied sentence-by-sentence between extracted reference sentences $\ell_r$
and corresponding planned outline sentences $\ell_d$ to produce the reconciled context $\mathcal{C}$.
\end{definition}

The design rationale for using LCS rather than cosine similarity over token embeddings
is threefold.
First, LCS is order-sensitive: it penalises rearrangements of key noun phrases,
which cosine similarity does not because bag-of-words embeddings treat a sentence as
a multiset of tokens, discarding order information.
For academic text in which the precise arrangement of technical terms encodes meaning
(e.g., ``neural network classification'' versus ``classification of neural networks''),
order sensitivity is a correctness requirement, not merely a preference.

Second, the threshold $\theta=4$ is directly interpretable by a domain expert:
it corresponds to a minimum of four shared tokens between the reference sentence and
the outline sentence, which in practice means that a key term or proper noun from
the source appears verbatim in the planned draft.
By contrast, a cosine threshold requires calibration against an embedding model and
has no direct linguistic interpretation.

Third, LCS is deterministic and parameter-free beyond $\theta$, making it auditable.
Given the same pair $(\ell_r,\ell_d)$, $\Phi_\theta$ always returns the same result.
This is a prerequisite for Invariant~I4 (Order Preservation), which requires the
reconciled context $\mathcal{C}$ to be a deterministic function of the outline $o$
and the reference set $\mathcal{R}^+$.

Figure~\ref{fig:lcs} illustrates the DP table for a representative sentence pair.
The highlighted diagonal cells trace the longest common subsequence; the dashed line
marks the $\theta=4$ acceptance threshold.

\begin{figure}[h!]
\centering
\begin{tikzpicture}[
  cell/.style={draw=DG!30, fill=white, minimum size=0.72cm,
               align=center, font=\scriptsize},
  hcell/.style={draw=C1, fill=L1, minimum size=0.72cm,
                align=center, font=\scriptsize\bfseries},
  hlbl/.style={font=\scriptsize\bfseries, text=C1},
  seq/.style={font=\scriptsize\bfseries, text=DG}
]
% Row/col labels — reference sequence ℓ_r
\foreach \i/\ch in {1/\varnothing,2/I,3/n,4/t,5/r,6/o,7/d,8/u,9/c}{
  \node[seq] at (\i*0.72-0.36, 7.2) {$\ch$};
}
% Column labels — draft sequence ℓ_d
\foreach \j/\ch in {1/\varnothing,2/I,3/n,4/v,5/e,6/s,7/t,8/i,9/g}{
  \node[seq] at (-0.52, -\j*0.72+7.2+0.36) {$\ch$};
}

% DP table (9x9, first row/col = 0 boundary)
\foreach \i in {1,...,9}{
  \foreach \j in {1,...,9}{
    \pgfmathsetmacro{\val}{%
      (\i==1 || \j==1) ? 0 :
      ((\i==2 && \j==2) ? 1 :
      ((\i==3 && \j==3) ? 2 :
      ((\i==4 && \j==4) ? 2 :
      ((\i==5 && \j==5) ? 2 :
      ((\i==6 && \j==6) ? 2 :
      ((\i==7 && \j==7) ? 2 :
      ((\i==8 && \j==8) ? 2 :
      ((\i==9 && \j==9) ? 2 : ""))))))))) }
    \node[cell] at (\i*0.72-0.36, -\j*0.72+7.92) {};
  }
}
% Fill known LCS diagonal (matching I, n)
\node[hcell] at (2*0.72-0.36, -2*0.72+7.92) {1};
\node[hcell] at (3*0.72-0.36, -3*0.72+7.92) {2};
% Highlight θ=4 threshold boundary
\draw[C4, dashed, line width=1.4pt] (0.0,7.6)--(6.5,1.1);
\node[font=\tiny\bfseries, text=C4] at (6.8,0.8){$\theta\!=\!4$};

% Traceback arrow along diagonal
\draw[-{Stealth[length=6pt,width=4pt]}, C1, line width=1.4pt]
  (2*0.72-0.36, -2*0.72+7.92) -- (3*0.72-0.36, -3*0.72+7.92);
\node[font=\tiny, text=C1, anchor=west] at (3*0.72, -3*0.72+8.05)
  {LCS traceback};

% Legend labels
\node[font=\scriptsize, text=MG, anchor=west] at (6.5,6.8)
  {$\ell_r\!=$ reference text};
\node[font=\scriptsize, text=MG, anchor=west] at (6.5,6.4)
  {$\ell_d\!=$ draft text};
\node[font=\scriptsize\bfseries, text=C1, anchor=west] at (6.5,6.0)
  {Matched cells};
\end{tikzpicture}
\caption{Schematic DP table for the LCS Reconciliation Operator
         $\Phi_\theta(\ell_r, \ell_d)$ (Definition~\ref{def:lcs}).
         Highlighted cells trace the longest common subsequence
         between an extracted reference segment and its draft
         counterpart. When $|\mathrm{LCS}(\ell_r,\ell_d)|\geq\theta=4$,
         the draft sentence is accepted without modification,
         preserving the author's phrasing while guaranteeing
         factual alignment with the source.}
\label{fig:lcs}
\end{figure}


\subsection{The Billing Operator}

\begin{definition}[Billing Operator $\mathcal{B}$]
\label{def:billing}
Let $T\in\mathbb{N}_0$ be the total input+output token count across all LLM calls in a
session, and $n_{\mathrm{img}}\in\mathbb{N}_0$ the number of TikZ figures generated.
The billing operator is the affine function
\[
  \mathcal{B}(T, n_{\mathrm{img}}) \;=\; 3T + 800\,n_{\mathrm{img}}
\]
measured in platform micro-units (PMU).
The cost is accumulated monotonically at each stage completion event.
\end{definition}

The linearity of $\mathcal{B}$ in both arguments is an intentional design constraint,
not a consequence of the underlying API pricing.
A non-linear billing function (e.g., tiered pricing) would complicate the pre-session
budget check: a user who enters with budget $\beta$ could not be guaranteed that spending
$\beta/2$ in the first half of the session leaves $\beta/2$ for the second half.
The affine structure of $\mathcal{B}$ guarantees this: if the first half of the pipeline
consumes $\mathcal{B}(T_1, n_1)$, the remaining budget is exactly
$\beta - \mathcal{B}(T_1, n_1)$.

The coefficient $800\,n_{\mathrm{img}}$ reflects the fixed overhead of invoking the
\LaTeX{} compiler as a correctness oracle for each figure: compilation, error parsing,
and up to $K=3$ repair iterations each consume a bounded number of tokens whose
expectation is approximately $266\,\text{PMU}$ per attempt (see Proposition~\ref{prop:token}).
Summing over three attempts gives the per-figure expected budget of
$3\times 266 \approx 800\,\text{PMU}$.

Figure~\ref{fig:billing} shows the iso-cost curves of $\mathcal{B}$ as a function of $T$
for five image counts, and marks a representative operating point.

\begin{figure}[h!]
\centering
\begin{tikzpicture}
\begin{axis}[PL,
  xlabel={Token count $T$ (thousands)},
  ylabel={Billed cost $\mathcal{B}(T, n_{\mathrm{img}})$ (micro-units)},
  title={Billing operator iso-curves for $n_{\mathrm{img}} \in \{0,1,2,4,8\}$},
  xmin=0, xmax=42, ymin=0, ymax=16000,
  xtick={0,5,10,15,20,25,30,35,40},
  ytick={0,2000,4000,6000,8000,10000,12000,14000,16000},
  legend pos=north west,
  width=13cm, height=7cm,
]
% B(T, n) = 3*T + 800*n  (T in thousands => multiply x by 1000)
\addplot[C1, very thick] coordinates{
  (0,0)(5,15000)(10,30000)(15,45000)(20,60000)(25,75000)(30,90000)(35,105000)(40,120000)};
% Scale: for readability show T in thousands but compute in hundreds
% B(T,0)=3T, B(T,1)=3T+800, etc. — keep raw T, annotate lines
\addplot[C1, very thick, domain=0:40]{3*x*1000};
\addlegendentry{$n_{\mathrm{img}}=0$}
\addplot[C2, thick, domain=0:40]{3*x*1000 + 800};
\addlegendentry{$n_{\mathrm{img}}=1$}
\addplot[C3, thick, domain=0:40]{3*x*1000 + 1600};
\addlegendentry{$n_{\mathrm{img}}=2$}
\addplot[C4, thick, domain=0:40]{3*x*1000 + 3200};
\addlegendentry{$n_{\mathrm{img}}=4$}
\addplot[C5, thick, domain=0:40]{3*x*1000 + 6400};
\addlegendentry{$n_{\mathrm{img}}=8$}
% Highlight typical operating point: T~12k tokens, n_img=3
\addplot[C4, mark=*, only marks, mark size=5pt] coordinates{(12,{3*12000+2400})};
\node[font=\tiny\bfseries, text=C4, anchor=west]
  at (axis cs:12.4,{3*12000+2400}){Typical\\session};
% Budget ceiling line
\addplot[DG, dashed, thin, domain=0:40]{10000};
\node[font=\tiny, text=DG, anchor=west] at (axis cs:30,10400){Budget cap};
\end{axis}
\end{tikzpicture}
\caption{Iso-cost curves of the Billing Operator $\mathcal{B}(T, n_{\mathrm{img}})=3T+800\,n_{\mathrm{img}}$
         (Definition~\ref{def:billing}) for five image-count values.
         The image penalty term dominates at high $n_{\mathrm{img}}$,
         establishing that visual richness is the primary cost driver.
         The annotated operating point represents a typical platform session.}
\label{fig:billing}
\end{figure}


\subsection{The Reference Quality Poset}

\begin{definition}[Reference Quality Poset $\mathcal{L}_Q$]
\label{def:quality_poset}
Define a partial order $\prec$ on the set of reference quality tiers
$\{Q_0,Q_1,Q_2,Q_3,Q_4,Q_5\}$ by:
\[
  Q_0 \prec Q_1 \prec Q_2 \prec Q_3 \prec Q_4 \prec Q_5
\]
where the tiers are defined as:
$Q_0$ = unresolvable or hallucinated (no live endpoint);
$Q_1$ = live web resource (HTTPS, last-modified header verifiable);
$Q_2$ = book or thesis (ISBN or institutional handle resolvable);
$Q_3$ = preprint with citation count $\geq\tau_s$;
$Q_4$ = conference proceedings indexed in DBLP or ACM DL;
$Q_5$ = peer-reviewed journal article with verified DOI.
The tuple $(\{Q_0,\ldots,Q_5\},\prec)$ is a totally ordered set (chain), hence a lattice.
\end{definition}

The poset $\mathcal{L}_Q$ serves two purposes in the system.
First, it provides a vocabulary for the \cgrag{} gate: a reference receives the
\texttt{enhanced} classification only if its tier is at least $Q_2$, which corresponds
to a verifiable academic source.
Second, it provides a partial order for the quality sub-metric $R(s)$: the quality
of the reference set $\mathcal{R}^+$ is the weighted mean of tier values, normalised
to $[0,1]$, providing a scalar quality signal that is monotone in the quality of
individual references.

The choice to define $\mathcal{L}_Q$ as a total order (chain) rather than a more
complex partial order reflects a deliberate simplicity trade-off.
One could argue that a preprint with 1,000 citations ($Q_3$) should be ranked above
a conference paper with 5 citations ($Q_4$), and that a partial order would capture this.
However, a total order enables the quality metric $R(s)$ to be computed as a simple mean
without requiring weighted aggregation of incomparable elements, and is therefore more
auditable and explainable to end users.

Figure~\ref{fig:hasse} shows the Hasse diagram of $\mathcal{L}_Q$ with the \cgrag{}
gate boundary overlaid.

\begin{figure}[h!]
\centering
\begin{tikzpicture}[
  tier/.style={draw=#1, fill=#2, rounded corners=4pt,
               minimum width=5.8cm, minimum height=0.80cm,
               align=center, font=\small\bfseries, line width=1.5pt},
  arr/.style={-{Stealth[length=7pt,width=5pt]}, line width=1.4pt, draw=DG!60},
  lbl/.style={font=\scriptsize, text=MG, anchor=west}
]
\node[tier={C1}{L1}] (Q5) at (0,5.6) {$Q_5$: Peer-reviewed journal (DOI verified)};
\node[tier={C2}{L2}] (Q4) at (0,4.2) {$Q_4$: Conference proceedings (DBLP indexed)};
\node[tier={C3}{L3}] (Q3) at (0,2.8) {$Q_3$: Preprint + citation count $\geq\tau_s$};
\node[tier={C4}{L4}] (Q2) at (0,1.4) {$Q_2$: Book / thesis (ISBN / handle resolved)};
\node[tier={C5}{L5}] (Q1) at (0,0.0) {$Q_1$: Web resource (HTTPS, last-modified)};
\node[tier={MG}{LGr}] (Q0) at (0,-1.4){$Q_0$: Unresolvable / hallucinated reference};

\draw[arr] (Q0)--(Q1); \node[lbl] at (0.15,-0.7){$\prec$};
\draw[arr] (Q1)--(Q2); \node[lbl] at (0.15, 0.7){$\prec$};
\draw[arr] (Q2)--(Q3); \node[lbl] at (0.15, 2.1){$\prec$};
\draw[arr] (Q3)--(Q4); \node[lbl] at (0.15, 3.5){$\prec$};
\draw[arr] (Q4)--(Q5); \node[lbl] at (0.15, 4.9){$\prec$};

% CGRAG gate indicator
\draw[C4, dashed, line width=1.2pt] (-3.1,1.0)--(3.1,1.0);
\node[font=\scriptsize\bfseries, text=C4] at (3.5,1.0){$\tau_s\!=\!100$};
\node[font=\scriptsize, text=C4] at (-4.1,1.3){\cgrag{} gate};
\node[font=\scriptsize\itshape, text=C3] at (-4.1,3.2){enhanced};
\node[font=\scriptsize\itshape, text=C4] at (-4.1,0.3){degraded};
\end{tikzpicture}
\caption{Hasse diagram of the Reference Quality Poset $\mathcal{L}_Q$
         (Definition~\ref{def:quality_poset}).
         Arrows denote the strict partial order $\prec$; each tier
         subsumes verifiability and academic authority criteria.
         The dashed line marks the \cgrag{} gate boundary at citation
         threshold $\tau_s\!=\!100$: references at $Q_2$ and above
         receive the \texttt{enhanced} classification.}
\label{fig:hasse}
\end{figure}


\subsection{The Multi-Criteria Lattice}

\begin{definition}[Multi-Criteria Lattice for \cgrag{}]
\label{def:lattice}
The three binary admission predicates
$p_c = [c\geq\tau_c]$, $p_s = [s\geq\tau_s]$, $p_k = [k\geq\tau_k]$
generate the Boolean lattice $2^{\{c,s,k\}}$ ordered by subset inclusion.
The meet operation $\wedge$ corresponds to logical conjunction of predicates,
and the join $\vee$ to logical disjunction.
The gate $\Gamma$ is the unique meet-semilattice homomorphism
$\Gamma: 2^{\{c,s,k\}} \to \{\texttt{degraded},\texttt{enhanced}\}$
that maps only the top element $\{c,s,k\}$ to \texttt{enhanced}.
\end{definition}

Definition~\ref{def:lattice} extends Definition~\ref{def:cgrag} by making the
algebraic structure of $\Gamma$ explicit.
The Boolean lattice $2^{\{c,s,k\}}$ has eight elements (the eight subsets of
$\{c,s,k\}$), arranged in four levels (Figure~\ref{fig:lattice}).
The homomorphism $\Gamma$ collapses all seven non-top elements to \texttt{degraded}
and maps the top element $\{c,s,k\}$ to \texttt{enhanced}.

The meet-semilattice structure is significant because it rules out the possibility of
\emph{compensatory} admission: a reference with a very high confidence score ($c\gg\tau_c$)
cannot be admitted despite failing the citation-count criterion ($s<\tau_s$).
In compensatory scoring systems (such as weighted sums), a very high score on one criterion
can outweigh a low score on another.
The lattice formulation explicitly rejects this: the gate is a conjunction, not a sum.

\subsection{Mutual Information for Context Relevance}

\begin{definition}[Mutual Information for Context--Section Relevance]
\label{def:mutual_info}
Let $X$ be the random variable representing the topic distribution of a retrieved
reference and $Y$ the random variable representing the target section distribution.
Both are modelled as categorical distributions over a shared topic vocabulary
$\mathcal{V}$ of size $|\mathcal{V}|=512$ topics extracted by LDA.
The context--section relevance score is
\[
  I(X;Y) \;=\; \sum_{x,y} P(X=x,Y=y)\log\frac{P(X=x,Y=y)}{P(X=x)P(Y=y)}.
\]
A reference is retained in the context window of $F_3$ for section $s_j$ only if
$I(X;Y)\geq I_{\min}$ for a session-configured threshold $I_{\min}$.
\end{definition}

This definition formalises the intuition that a reference about quantum computing
should not appear in the context window of a section about database indexing, even if
both achieve \texttt{enhanced} classification under $\Gamma$.
The mutual information criterion is complementary to $\Gamma$: $\Gamma$ filters on
scholarly authority, and $I(X;Y)\geq I_{\min}$ filters on topical relevance.

The use of LDA topics as the shared vocabulary $\mathcal{V}$ is motivated by the
need for a topic representation that is independent of the LLM generating the draft.
Using the LLM's own latent representations to assess relevance would risk circular
reasoning: a reference that the LLM has internally associated with the section topic
would be selected, reinforcing whatever biases the LLM already has about that topic.
LDA topics, estimated from the reference corpus independently of the LLM, provide
an external relevance signal.

\subsection{The Typed Visual Grammar}

\begin{definition}[Typed Visual Grammar \tvg{}]
\label{def:tvg}
A Typed Visual Grammar is a 4-tuple $G=(T,B,L,Q)$ where:
$T=\{t_1,\ldots,t_{30}\}$ is a finite set of figure types
(e.g., \texttt{flow\_diagram}, \texttt{state\_machine}, \texttt{bar\_chart},
\texttt{scatter\_plot}, \texttt{heatmap}, \texttt{hierarchy}, \texttt{timeline});
$B: T\to\Sigma^*$ is a blueprint function mapping each type to a \LaTeX{}/TikZ
structural template;
$L$ is a \LaTeX{} compiler invoked as a first-class correctness oracle;
$Q: \Sigma^*\times\Sigma^*\to\{\mathrm{accept},\mathrm{repair}(e)\}$
is the quality judgment that maps (tikzCode, stderr) to either acceptance or a
structured repair directive $e$.
The \tvg{} cycle is: generate $\to$ compile via $L$ $\to$ $Q$ classifies
$\to$ if repair, feed $e$ back to generation; repeat for at most $K=3$ attempts.
\end{definition}

The type set $T$ with $|T|=30$ is designed to cover all major figure categories
that appear in academic publications across computer science, social science,
biomedical informatics, and engineering.
The 30 types include: 7 flow/process types (linear flow, forked flow, cyclic flow,
swim-lane, pipeline, state machine, event sequence); 6 quantitative types
(bar chart, line chart, scatter plot, histogram, heatmap, box plot);
5 structural types (tree, hierarchy, mind map, matrix, table);
4 architectural types (layer diagram, network graph, component diagram, deployment map);
4 relational types (Venn diagram, lattice, Hasse diagram, bipartite graph);
and 4 temporal types (timeline, Gantt chart, calendar, phase diagram).

The blueprint function $B: T\to\Sigma^*$ is the intellectual-property core of the
\tvg{} system.
Each blueprint is a TikZ/pgfplots template with named placeholder variables
(figure title, axis labels, data series, node labels) that the LLM fills in from
the section context.
The blueprint constrains the \emph{structure} of the generated code while leaving
the \emph{content} to be determined by the LLM.
This separation of structure and content is what makes the compile-repair cycle
tractable: structural errors (missing \texttt{\textbackslash end\{tikzpicture\}},
undefined TikZ styles) are rare because the blueprint enforces the overall structure,
and the remaining errors are content errors (incorrect numeric values, invalid colour
names) that the LLM can fix from the stderr message.

The key design decision is the use of $L$ (the \LaTeX{} compiler) as an oracle
rather than as a post-hoc validator.
By invoking the compiler \emph{during} the generation loop, the system obtains two
pieces of information with each compilation attempt: a binary correctness signal
(exit code 0/non-zero) and a structured error signal (stderr) that is directly
interpretable by the LLM in the next repair iteration.
This is in contrast to approaches that validate figures only at document-assembly
time (Stage~$F_4$), where a compilation error is unrecoverable because the assembler
has already embedded the figure \texttt{\textbackslash input} directive in the body.
The compile-validate-repair cycle is illustrated in Figure~\ref{fig:tikz_repair}.

\begin{figure}[h!]
\centering
\begin{tikzpicture}[
  circ/.style={draw=#1, fill=#2, circle, minimum size=2.0cm,
               align=center, font=\scriptsize\bfseries, inner sep=2pt,
               line width=1.5pt},
  rect/.style={draw=#1, fill=#2, rounded corners=3pt,
               minimum width=2.8cm, minimum height=0.75cm,
               align=center, font=\scriptsize\bfseries, line width=1.2pt},
  arr/.style={-{Stealth[length=7pt,width=5pt]}, line width=1.4pt, draw=#1},
  lbl/.style={font=\tiny, text=DG, align=center, text width=2.4cm}
]
\node[rect={DG}{LGr}] (bp) at (0,0) {30-Type\\Blueprint Grammar};
\node[circ={C1}{F1}]  (gen)  at (90:3.4)  {LLM\\Generates\\TikZ};
\node[circ={C2}{F2}]  (comp) at (0:3.4)   {Go\\Compiler\\(standalone)};
\node[circ={C3}{F3}]  (ok)   at (270:3.4) {Accept\\+\\Embed};
\node[circ={C4}{F4}]  (rep)  at (180:3.4) {Feed\\stderr\\to LLM};

\draw[arr={C1}] (bp)  to[bend left=10]  (gen);
\draw[arr={DG}] (gen) to[bend left=10]  node[lbl, right=2pt]{tikzCode} (comp);
\draw[arr={C3}] (comp) to[bend left=10] node[lbl, below=2pt]{OK} (ok);
\draw[arr={C4}] (comp) to[bend left=10] node[lbl, left=2pt]{FAIL (stderr)} (rep);
\draw[arr={DG}] (rep) to[bend left=10]  node[lbl, above=2pt]{attempt $\leq3$} (gen);

\node[rect={LGr}{LGr}, minimum width=2.4cm] (out) at (270:5.0)
  {To Stage 4 Assembler};
\node[rect={C4}{L4}, minimum width=2.2cm] (fail) at (235:5.2)
  {Mark failed\\(skip figure)};
\draw[arr={C3}] (ok)--(out);
\draw[arr={C4}, dashed] (rep) to[bend right=18] (fail);
\end{tikzpicture}
\caption{Compile-validate-repair cycle for the Typed Visual Grammar system
         (\tvg{}, Definition~\ref{def:tvg}). The \LaTeX{} compiler functions
         as a formal correctness oracle: its exit code constitutes proof of
         structural validity; its \texttt{stderr} provides a structured
         error signal for the next repair iteration.}
\label{fig:tikz_repair}
\end{figure}


\subsection{The Language Partition}

\begin{definition}[Language Partition $\mathcal{L}_{25}$]
\label{def:language}
Define four mutually exclusive language classes:
$L_{\mathrm{kk}}$ = Kazakh (Cyrillic script, ISO 639-1 \texttt{kk}, as used in
Kazakh academic publishing prior to the Latin transition);
$L_{\mathrm{cyr}}$ = generic Cyrillic academic languages (Russian, Belarusian, Uzbek,
and other CIS languages using the Cyrillic alphabet);
$L_{\mathrm{lat}}$ = Kazakh (Latin script, as standardised by the 2021 Kazakh Latin
alphabet adopted by the Kazakhstani government);
$L_{\mathrm{en}}$ = English and other Latin-script academic languages (German, French,
Spanish) used in international venues.
The partition is:
\[
  \mathcal{L}_{25} \;=\; L_{\mathrm{kk}} \sqcup L_{\mathrm{cyr}} \sqcup L_{\mathrm{lat}} \sqcup L_{\mathrm{en}}
\]
where $\sqcup$ denotes disjoint union.
Every session $\sigma$ is assigned to exactly one class $\ell(\sigma)\in\mathcal{L}_{25}$
by a deterministic language-detection step at the start of $F_1$.
\end{definition}

The partition is exhaustive: since $\mathcal{L}_{25}$ covers all scripts present
in Kazakh academic publishing (the primary deployment context), no session can fail
to be assigned a language class.
Exhaustiveness is enforced structurally: the language-detection step returns a sum type
with four variants corresponding to the four classes, and the subsequent prompting
logic pattern-matches exhaustively on all four variants with no default or fallthrough case.

The name $\mathcal{L}_{25}$ reflects the platform's original specification year
(2025) and is not related to a cardinality of 25 languages.
The subscript disambiguates this object from other uses of the letter $L$ in the
formal model (lattice, LCS).

An important design constraint of $\mathcal{L}_{25}$ is the strict separation of
$L_{\mathrm{kk}}$ (Kazakh Cyrillic) from $L_{\mathrm{cyr}}$ (generic Cyrillic).
In many multilingual NLP systems, Kazakh Cyrillic and Russian Cyrillic are conflated
because they share the same Unicode code points.
\sys{} separates them because the prompting conventions, citation styles (GOST vs.
APA), and keyboard encoding rules differ between the two, and a system that treats
them identically produces suboptimal output for Kazakh-language sessions.

\subsection{The Unified System Model}

\begin{definition}[Unified System Model $\mathcal{P}$]
\label{def:unified}
The \sys{} platform is the 8-tuple
\[
  \mathcal{P} \;=\; (G,\,\Gamma,\,\mathcal{B},\,M,\,\Phi_\theta,\,\mathcal{L}_Q,\,\mathcal{L}_{25},\,\Pi)
\]
where $G$ is the Typed Visual Grammar (Def.~\ref{def:tvg}),
$\Gamma$ is the \cgrag{} gate (Def.~\ref{def:cgrag}),
$\mathcal{B}$ is the Billing Operator (Def.~\ref{def:billing}),
$M$ is the session metadata record containing (user ID, topic string, budget ceiling
$\beta$, session timestamp $t_0$),
$\Phi_\theta$ is the LCS Reconciliation Operator (Def.~\ref{def:lcs}),
$\mathcal{L}_Q$ is the Reference Quality Poset (Def.~\ref{def:quality_poset}),
$\mathcal{L}_{25}$ is the Language Partition (Def.~\ref{def:language}), and
$\Pi$ is the Pipeline Algebra (Def.~\ref{def:pipeline}).
\end{definition}

The choice to bundle $M$ (session metadata) into $\mathcal{P}$ as an explicit component
reflects the observation that session metadata is not merely administrative but participates
in formal reasoning.
The session timestamp $t_0$ is used in Invariant~I5 (Real-Data Guarantee) to determine
which cached reference entries are stale.
The budget ceiling $\beta$ is used in $\mathcal{B}$ for pre-session feasibility checks.
By including $M$ in $\mathcal{P}$, we make these dependencies explicit and subject to
formal reasoning.

Figure~\ref{fig:unified} shows the structural decomposition of $\mathcal{P}$ with its
eight components and their inter-component contracts.
The figure highlights the central role of $\Pi$: all other components are either
invoked by a specific stage of $\Pi$ or constrain the behaviour of those stages via
the invariant envelope.

\begin{figure}[h!]
\centering
\begin{tikzpicture}[
  comp/.style={draw=#1, fill=#2, rounded corners=4pt,
               minimum width=3.2cm, minimum height=0.85cm,
               align=center, font=\scriptsize\bfseries, line width=1.4pt},
  hub/.style={draw=C1, fill=L1, rounded corners=6pt,
              minimum width=3.8cm, minimum height=1.1cm,
              align=center, font=\small\bfseries, line width=2pt},
  arr/.style={-{Stealth[length=6pt,width=4pt]}, line width=1.2pt, draw=#1},
]
% Central hub — Unified System Model P
\node[hub] (P) at (0,0) {$\mathcal{P}$\\Unified Model};

% Eight components arranged around hub
\node[comp={C2}{L2}] (G)   at (0,    2.6) {$G$: Typed Visual\\Grammar (\tvg{})};
\node[comp={C3}{L3}] (gam) at (3.0,  1.6) {$\Gamma$: \cgrag{}\\Gate};
\node[comp={C4}{L4}] (B)   at (3.8,  0.0) {$\mathcal{B}$: Billing\\Operator};
\node[comp={C5}{L5}] (M)   at (3.0, -1.6) {$M$: Language\\Partition $\mathcal{L}_{25}$};
\node[comp={C6}{L6}] (phi) at (0,   -2.6) {$\Phi_\theta$: LCS\\Reconciliation};
\node[comp={C7}{L7}] (LQ)  at (-3.0,-1.6) {$\mathcal{L}_Q$: Quality\\Poset};
\node[comp={C8}{L4}] (L25) at (-3.8, 0.0) {$\mathcal{L}_{25}$: Kazakh\\+ Cyrillic + Lat};
\node[comp={C1}{L1}] (Pi)  at (-3.0, 1.6) {$\Pi$: Pipeline\\Algebra};

% Arrows from hub outward (bidirectional for most)
\draw[arr={C2}] (P)--(G);
\draw[arr={C3}] (P)--(gam);
\draw[arr={C4}] (P)--(B);
\draw[arr={C5}] (P)--(M);
\draw[arr={C6}] (P)--(phi);
\draw[arr={C7}] (P)--(LQ);
\draw[arr={C8}] (P)--(L25);
\draw[arr={C1}] (P)--(Pi);

% Invariants label
\node[font=\tiny\itshape, text=MG, text width=3cm, align=center]
  at (0,-3.8){Governed by Invariants I1--I5};
\draw[DG!40, dashed, rounded corners=8pt, line width=1pt]
  (-5.2,-3.2) rectangle (5.2,3.5);
\node[font=\tiny\bfseries, text=MG] at (4.5,3.2){$\mathcal{P}$};
\end{tikzpicture}
\caption{Structural decomposition of the Unified System Model
         $\mathcal{P}=(G,\Gamma,\mathcal{B},M,\Phi_\theta,\mathcal{L}_Q,\mathcal{L}_{25},\Pi)$
         (Definition~\ref{def:unified}).
         The eight formal components are joined at the central
         coordination point $\mathcal{P}$; each edge represents
         a formal dependency or data-flow contract.
         The outer boundary denotes the invariant envelope
         (Invariants~I1--I5) that constrains all inter-component
         interactions at runtime.}
\label{fig:unified}
\end{figure}


\subsection{Physical Architecture}

The physical realisation of $\mathcal{P}$ follows a six-layer separation of concerns,
illustrated in Figure~\ref{fig:architecture}.
Each layer encapsulates a distinct concern and communicates with adjacent layers only
through formally typed interfaces.
The six layers, from lowest to highest, are:

\paragraph{Data Layer.}
PostgreSQL\,16 is the authoritative store for session records, billing events,
reference quality scores ($\mathcal{L}_Q$ tier assignments), and compiled document artefacts.
All writes to the billing table are append-only to preserve Invariant~I3 (Billing Monotonicity):
once a billing event is recorded, it cannot be modified or deleted.
MinIO provides S3-compatible object storage for compiled PDFs and intermediate
\LaTeX{} source bundles, enabling persistent access to prior session outputs.

\paragraph{Compute Layer.}
A stateless Go-based \texttt{research-api} microservice implements all six pipeline
stages ($F_1$--$F_5$, including $F_{2.5}$).
The Go runtime's goroutine model is the physical substrate for the
\texttt{concurrentMap} primitive proven in Theorem~\ref{thm:concurrent}.
The \LaTeX{} compiler oracle (used by $F_5$ and the \tvg{} repair cycle) is invoked
as a sandboxed subprocess with resource limits to prevent runaway compilation from
blocking the pipeline.
Python and R sub-processes are available for statistical computation (used by the
quality metric $Q$) and are invoked via the same subprocess interface.

\paragraph{Real-Time Layer.}
Yjs CRDT over \texttt{y-websocket} provides conflict-free concurrent editing of
document outlines during the $F_1$ planning stage, allowing multiple collaborators
to work on the same outline without coordination overhead.
A Telegram Bot integration delivers session completion and billing notifications,
enabling asynchronous workflow (a user can submit a session and receive the compiled
PDF URL without keeping a browser tab open).

\paragraph{AI Layer.}
All LLM calls are routed through a \texttt{chatCompletionLong} abstraction that
transparently handles token-window chunking, retry on API rate limits, and billing
accumulation.
The abstraction exposes a single function signature regardless of whether the
underlying call requires one API round-trip (short context) or multiple sequential
round-trips with context carry-over (long context that exceeds the model's context
window).
This transparency is necessary for $\mathcal{B}$ to accumulate costs correctly:
the billing accumulation logic is inside \texttt{chatCompletionLong} and therefore
cannot be bypassed by code that calls the LLM API directly.

\paragraph{API and Auth Layer.}
Next.js\,16 App Router provides JWT-authenticated REST endpoints for session
submission, status polling (\texttt{stage\_json} endpoint), and billing queries.
Rate limiting at this layer prevents any single user from consuming more than their
pre-approved budget share, independently of the per-session budget ceiling $\beta$
stored in $M$.
A real-time billing dashboard exposes $\mathcal{B}(T,n_{\mathrm{img}})$ values
computed from the current session state, allowing users to monitor cost accumulation
during long-running sessions.

\paragraph{Presentation Layer.}
A React\,19 single-page application provides three primary interfaces:
the Space Editor (collaborative outline authoring, backed by the Yjs layer),
the Analytics Dashboard (session cost visualisation, using the billing data from
the Data layer), and the PDF Viewer (compiled document preview, backed by MinIO
artefact retrieval).
The Presentation Layer has no direct access to the Data or Compute layers;
all communication is mediated by the API layer, ensuring that the six-layer
downward-communication constraint is preserved.

\begin{figure}[h!]
\centering
\begin{tikzpicture}[
  layer/.style={draw=DG, line width=1.5pt, rounded corners=5pt,
                minimum width=13cm, minimum height=0.9cm,
                align=center, font=\small\bfseries, text width=12.5cm},
  badge/.style={font=\scriptsize\itshape, text=MG, anchor=east}
]
\node[layer, fill=LGr!70]  (L1) at (0,0.00){Presentation: React\,19 · Space Editor · Analytics Dashboard · PDF Viewer};
\node[layer, fill=F1!80]   (L2) at (0,1.10){API \& Auth: Next.js\,16 App Router · JWT Sessions · Billing · Rate Limiting};
\node[layer, fill=F5!80]   (L3) at (0,2.20){AI Pipeline: Stages 1--5 · OpenAI GPT-5-mini · chatCompletionLong()};
\node[layer, fill=F2!80]   (L4) at (0,3.30){Real-time: Yjs CRDT · y-websocket · Telegram Bot (Node.js)};
\node[layer, fill=F3!80]   (L5) at (0,4.40){Compute: pdf-extractor · r-compiler · python-compiler · research-api (Go)};
\node[layer, fill=F4!80]   (L6) at (0,5.50){Data: PostgreSQL\,16 · Prisma ORM\,7 · MinIO S3-compatible Storage};

\node[badge] at (-6.65,0.00) {UI};
\node[badge] at (-6.65,1.10) {API};
\node[badge] at (-6.65,2.20) {AI};
\node[badge] at (-6.65,3.30) {RT};
\node[badge] at (-6.65,4.40) {Compute};
\node[badge] at (-6.65,5.50) {Data};
\end{tikzpicture}
\caption{Six-layer platform architecture of \sys{}. Each layer isolates a
         distinct concern: data persistence, distributed compute, AI pipeline,
         real-time collaboration, API routing, and presentation. All inter-layer
         communication is strictly downward.}
\label{fig:architecture}
\end{figure}


\subsection{System Invariants}

The five system invariants (I1--I5) shown in Figure~\ref{fig:invariants} are properties
that hold for every session $\sigma\in\mathcal{P}$ by construction.
We state each invariant formally and then sketch the proof of its constructive validity.

\begin{figure}[h!]
\centering
\begin{tikzpicture}[
  inv/.style={draw=#1, fill=#2, rounded corners=5pt,
              minimum width=13.5cm, minimum height=0.95cm,
              align=left, font=\scriptsize, line width=1.5pt,
              text width=13.0cm, inner sep=6pt},
  idlbl/.style={font=\scriptsize\bfseries, text=#1}
]
\node[inv={C1}{L1}] (I1) at (0, 4.8)
  {\textbf{I1 — Preamble Uniqueness.}
   $\forall d\in\mathcal{D}:\,|\{\text{preamble blocks in }d\}|=1$.
   Every generated document contains exactly one \LaTeX{} preamble;
   duplicate package declarations are structurally impossible.};

\node[inv={C2}{L2}] (I2) at (0, 3.3)
  {\textbf{I2 — Citation Safety.}
   $\forall r\in\text{refs}(d):\,\mathcal{Q}(r)\geq Q_1$.
   No reference with $Q_0$ (unresolvable) status survives into the
   assembled document; the \cgrag{} gate removes all sub-threshold entries.};

\node[inv={C3}{L3}] (I3) at (0, 1.8)
  {\textbf{I3 — Billing Monotonicity.}
   $\mathcal{B}(T_1,n_1)\leq\mathcal{B}(T_2,n_2)$ whenever $T_1\!\leq\!T_2$
   and $n_1\!\leq\!n_2$.
   Cost is strictly non-decreasing in both token usage and image count;
   refunds are architecturally impossible.};

\node[inv={C4}{L4}] (I4) at (0, 0.3)
  {\textbf{I4 — Order Preservation.}
   $\forall i<j:\,\text{section}_i$ appears before $\text{section}_j$
   in the compiled document. The \texttt{concurrentMap} primitive
   (Theorem~\ref{thm:concurrent}) guarantees index-stable emission
   despite parallel execution.};

\node[inv={C5}{L5}] (I5) at (0,-1.2)
  {\textbf{I5 — Real-Data Guarantee.}
   $\forall r\in\text{refs}(d):\,r$ was fetched from a live HTTP
   endpoint at generation time.
   No reference is fabricated or retrieved from a stale cache older
   than the session timestamp $t_0$.};

% Brace labelling
\draw[decorate, decoration={brace, amplitude=8pt, raise=2pt},
      C1, line width=1.2pt]
  (7.1,5.35) -- (7.1,-1.75)
  node[midway, right=10pt, font=\scriptsize\bfseries, text=C1]
    {$\mathcal{P}$ invariants};
\end{tikzpicture}
\caption{The five system invariants (I1--I5) of the \sys{} Unified Model $\mathcal{P}$.
         Each invariant is a universally quantified property
         that holds across all platform sessions by construction —
         not as a runtime check but as a consequence of the
         formal design decisions encoded in Definitions~\ref{def:pipeline}
         through~\ref{def:unified}.}
\label{fig:invariants}
\end{figure}


\begin{invariant}[Preamble Uniqueness]
\label{inv:preamble}
$\forall\,d\in\mathcal{D}: |\{\text{preamble blocks in }d\}|=1$.
\end{invariant}

\noindent\textit{Proof sketch.}
Stage $F_4$ (Assembler) inserts exactly one preamble block — the one imported from
the session preamble template — at the beginning of the assembled document $d_0$.
No other stage inserts preamble material.
Since stage execution is sequential and $F_4$ is the only stage that writes the
document header, the invariant holds. \hfill$\square$

\begin{invariant}[Citation Safety]
\label{inv:citation}
$\forall\,r\in\mathrm{refs}(d): \mathcal{Q}(r)\geq Q_1$.
\end{invariant}

\noindent\textit{Proof sketch.}
Stage $F_2$ produces $\mathcal{R}^+$ by filtering all candidates through $\Gamma$.
References with $\Gamma=\texttt{degraded}$ — which includes all $Q_0$ references — are
excluded from $\mathcal{R}^+$.
Stage $F_4$ assembles the bibliography from $\mathcal{R}^+$ exclusively.
Therefore, no $Q_0$ reference can appear in $\mathrm{refs}(d)$. \hfill$\square$

\begin{invariant}[Billing Monotonicity]
\label{inv:billing}
$T_1\leq T_2\wedge n_1\leq n_2\Rightarrow\mathcal{B}(T_1,n_1)\leq\mathcal{B}(T_2,n_2)$.
\end{invariant}

\noindent\textit{Proof sketch.}
$\mathcal{B}(T,n)=3T+800n$ is linear with non-negative coefficients $3$ and $800$.
Linear functions with non-negative coefficients are isotone (order-preserving) in
each argument.
Therefore, $T_1\leq T_2$ implies $3T_1\leq 3T_2$, and $n_1\leq n_2$ implies
$800n_1\leq 800n_2$; summing gives the result. \hfill$\square$

\begin{invariant}[Order Preservation]
\label{inv:order}
$\forall\,i<j: \mathrm{section}_i$ appears before $\mathrm{section}_j$ in the
compiled document output of $F_5$.
\end{invariant}

\noindent\textit{Proof sketch.}
Follows from Theorem~\ref{thm:concurrent}(a): the \texttt{concurrentMap} primitive
preserves the index order of its output regardless of worker completion order.
Stage $F_4$ (Assembler) reads the \texttt{concurrentMap} output in index order and
concatenates the sections in that order. \hfill$\square$

\begin{invariant}[Real-Data Guarantee]
\label{inv:realdata}
$\forall\,r\in\mathrm{refs}(d)$: $r$ was fetched from a live HTTP endpoint at
generation time $t_0$ of session $\sigma$; no reference is retrieved from a
stale cache older than $t_0$.
\end{invariant}

\noindent\textit{Proof sketch.}
The \cgrag{} retrieval protocol issues an HTTP request with
\texttt{If-Modified-Since:}~$t_0$ for every candidate reference.
A cached response with \texttt{Last-Modified} earlier than $t_0$ is rejected and
a fresh fetch is issued.
Therefore, all admitted references have been live-verified at or after $t_0$.
This proof is conditional on network liveness at $t_0$. \hfill$\square$
